% !TEX root = ../TRENCH_DIFFEQ.tex

\chapter{Applications of Linear Second Order Equations}
\label{ch:apps_second}

\noindent
\textsc{In this chapter}
 we study applications of linear second order
equations.

\medskip\noindent
\zcref[S,CAP,font=\scshape]{sec:spring_i,sec:spring_ii} are about spring-mass systems.



\medskip\noindent
\zcref[S,CAP,font=\scshape]{sec:rlc} is about $RLC$ circuits, the electrical
analogs of spring-mass systems.


\medskip\noindent
\zcref[S,CAP,font=\scshape]{sec:central_force} is  about
motion of an object under a central force, which is
particularly relevant in the space age, since, for example, a
satellite moving in orbit subject only to Earth's gravity is
experiencing motion under a central force.

\newpage

\section{Spring Problems I}\label{sec:spring_i}
\index{spring problems}

We consider the motion of an object of mass $m$, suspended from a
spring of negligible mass. We say that the spring-mass system is
in \emph{\color{blue}equilibrium}\index{equilibrium!spring-mass system}
when the object is at rest and the forces
acting on it sum to zero. The position of the object in this case is
the \indexterm{equilibrium position}. We define $y$ to be the displacement
of the object from its equilibrium position (\zcref{figure:6.1.1}),
measured positive upward.


\noindent
\begin{minipage}[t]{0.48\linewidth}
\tagstructbegin{tag=Sect}
\includegraphics[alt={Three springs with masses that are (a) above, (b) at, or (c) below the equilibrium y=0.},
width=\linewidth,height=2.44in,keepaspectratio]{fig060101}
\captionof{figure}{(a) $y>0$, (b) $y=0$, (c) $y<0$}
\label{figure:6.1.1}
\tagstructend
\end{minipage}%
\hfill
\begin{minipage}[t]{0.48\linewidth}
\tagstructbegin{tag=Sect}
\includegraphics[alt={A spring with a mass connected to a U shaped damper.},
width=\linewidth,height=2.44in,keepaspectratio]{fig060102}
\captionof{figure}{\raggedright A spring-mass system with damping}
\label{figure:6.1.2}
\tagstructend
\end{minipage}


Our model accounts for the following kinds of forces acting on the
object:

\begin{itemize}
\item The force $-mg$, due to gravity.

\item A force $F_s$ exerted by the spring resisting change in its
length. The \emph{\color{blue}natural length} of the spring is its length with
no mass attached. We assume that the spring obeys
\personHref{Hooke}{Hooke's law}\index{Hooke's law}:
If the length of the spring is changed by an amount $\Delta L$
from its natural length,  then the spring exerts a force $F_s=k\Delta
L$, where $k$ is a positive number called the \emph{\color{blue}spring constant}.
If the spring is stretched then  $\Delta L>0$ and $F_s>0$, so the
spring force is upward, while if the spring is compressed then $\Delta
L<0$ and $F_s<0$, so the spring force is downward.

\item A \emph{\color{blue}damping force} $F_d=-cy'$ that resists the motion with
a force proportional to the velocity of the object. It may be due to
air resistance or friction in the spring. However, a convenient way to
visualize a damping force is to assume that the object is rigidly
attached to a piston with negligible mass immersed in a cylinder
(called a \indexterm{dashpot}) filled with a viscous liquid
(\zcref{figure:6.1.2}). As the piston moves, the liquid exerts a
damping force. We say that the motion is \emph{\color{blue}undamped} if $c=0$, or
\emph{\color{blue}damped} if $c>0$.


\item An external force $F$, other than the force due to gravity, that
may vary with $t$, but is independent of displacement and velocity. We
say that the motion is \emph{\color{blue}free} if $F\equiv0$, or \emph{\color{blue}forced}
if $F\not\equiv0$.

\end{itemize}



From Newton's second law of motion\index{Newton's!second law of motion},
\begin{equation}\label{eq:6.1.1}
my''=-mg+F_d+F_s+F=-mg-cy'+F_s+F.
\end{equation}
We must now relate $F_s$ to $y$. In the absence of external forces the
object stretches the spring by an amount $\Delta l$ to assume its
equilibrium position (\zcref{figure:6.1.3}). Since the sum of the forces acting on the object
is then zero, Hooke's Law implies that $mg=k\Delta l$. If the object
is displaced $y$ units from its equilibrium position, the total
change in the length of the spring is $\Delta L=\Delta l-y$,
 so Hooke's law implies that
 \[
F_s=k\Delta L=k\Delta l-ky.
\]
Substituting this into \eqref{eq:6.1.1} yields
\[
my''=-mg-cy'+k\Delta L-ky+F.
\]
Since $mg=k\Delta l$ this can be written as
\begin{equation}\label{eq:6.1.2}
my''+cy'+ky=F.
\end{equation}
We call this \emph{\color{blue}the equation of motion}.

\begin{figure}
  \centering
  \includegraphics[alt={(a) A spring of length L.  (b) The same spring, stretched by a mass a further length ∆L.},
  width=5.67in,height=3.3in,keepaspectratio]{fig060103}
\color{blue}
\caption{(a) Natural length of spring, (b) Spring stretched by mass}
  \label{figure:6.1.3}
\end{figure}


\boxit{Simple Harmonic Motion}
\index{simple harmonic motion}

Throughout the rest of this section we'll consider spring-mass
systems without damping; that is, $c=0$. We'll consider systems with
damping in the next section.

We first consider the case where the motion is also free; that is,
$F$=0. We begin with an example.

\begin{example}\label{example:6.1.1}
An object stretches a spring 6 inches in equilibrium.
\begin{alist}
\item % (a)
Set up the equation of motion and find its general solution.

\item % (b)
Find the displacement of the object for $t>0$ if it's initially
displaced 18 inches above equilibrium and given a downward velocity of
\qty{3}{\feet\per\second}.
\end{alist}
\end{example}

\begin{alistSolution}
\item
Setting $c=0$ and $F=0$ in \eqref{eq:6.1.2} yields the equation of motion
\[
my''+ky=0,
\]
which we rewrite as
\begin{equation}\label{eq:6.1.3}
y''+\frac{k}{m}y=0.
\end{equation}
Although we would need the weight of the object to obtain $k$ from the
equation $mg=k\Delta l$ we can obtain $k/m$ from $\Delta l$ alone;
thus, $k/m=g/\Delta l$. Consistent with the units used in the problem
statement, we take $g=\qty{32}{\feet\per\second\squared}$. Although $\Delta l$ is stated in
inches, we must convert it to feet to be consistent with this choice
of $g$; that is, $\Delta l =1/2$ \unit{\feet}. Therefore
\[
\frac{k}{m}=\frac{32}{1/2}=64
\]
and \eqref{eq:6.1.3} becomes
\begin{equation}\label{eq:6.1.4}
y''+64y=0.
\end{equation}
 The characteristic equation of \eqref{eq:6.1.4} is
\[
r^2+64=0,
\]
which has the zeros $r=\pm 8i$. Therefore the general solution of
\eqref{eq:6.1.4} is
\begin{equation}\label{eq:6.1.5}
y=c_1\cos8t+c_2\sin8t.
\end{equation}

\item
The initial upward displacement of 18 inches is positive and must be
expressed in feet. The initial downward velocity is negative;   thus,
\[
y(0)=\frac{3}{2}\quad\text{and}\quad  y'(0)=-3.
\]
 Differentiating \eqref{eq:6.1.5} yields
\begin{equation}\label{eq:6.1.6}
y'=-8c_1\sin8t+8c_2\cos8t.
\end{equation}
Setting $t=0$ in \eqref{eq:6.1.5} and \eqref{eq:6.1.6} and imposing the
initial conditions shows that $c_1=3/2$ and $c_2=-3/8$. Therefore
\[
y=\frac{3}{2}\cos8t-\frac{3}{8}\sin8t,
\]
where $y$ is in feet (\zcref{figure:6.1.4}).

\begin{figure}
  \centering
  \includegraphics[alt={A sinusoidal curve, with amplitude about 1.5 and period about 0.8},
  width=5.67in,height=3.3in,keepaspectratio]{fig060104}
\color{blue}
\caption{$y=\dfrac{3}{2}\cos8t-\dfrac{3}{8}\sin 8t$}
  \label{figure:6.1.4}
\end{figure}
\end{alistSolution}

We'll now consider the equation
\[
my''+ky=0
\]
where $m$ and $k$ are arbitrary positive numbers. Dividing through by
$m$ and defining $\omega_0=\sqrt{k/m}$ yields
\[
y''+\omega_0^2y=0.
\]
The general solution of this equation is
\begin{equation}\label{eq:6.1.7}
y=c_1\cos\omega_0t+c_2\sin\omega_0t.
\end{equation}
We can rewrite this in a more useful form by defining
\begin{equation}\label{eq:6.1.8}
R=\sqrt{c_1^2+c_2^2},
\end{equation}
and
\begin{equation}\label{eq:6.1.9}
c_1=R\cos\theta\quad\text{and}\quad c_2=R\sin\theta.
\end{equation}
Substituting from \eqref{eq:6.1.9} into \eqref{eq:6.1.7} and applying the
identity
\[
\cos\omega_0t\cos\theta+\sin\omega_0t\sin\theta=\cos(\omega_0t-\theta)
\]
yields
\begin{equation}\label{eq:6.1.10}
y=R\cos(\omega_0t-\theta).
\end{equation}

From \eqref{eq:6.1.8} and \eqref{eq:6.1.9} we see that the $R$ and $\phi$
can be interpreted as polar coordinates of the point with rectangular
coordinates $(c_1,c_2)$ (\zcref{figure:6.1.5}). Given $c_1$ and
$c_2$, we can compute $R$ from \eqref{eq:6.1.8}. From \eqref{eq:6.1.8} and
\eqref{eq:6.1.9}, we see that $\theta$ is related to $c_1$ and $c_2$ by
\[
\cos\theta=\frac{c_1}{\sqrt{c_1^2+c_2^2}}\quad\text{and}\quad
\sin\theta=\frac{c_2}{\sqrt{c_1^2+c_2^2}}.
\]
There are infinitely many angles $\theta$, differing
 by integer multiples of $2\pi$, that satisfy these equations. We
will always choose $\theta$ so that $-\pi\le\theta<\pi$.

\begin{figure}
  \centering
  \includegraphics[alt={A point in the first quadrant with coordinates (c1,c2) that is distance R from the origin and makes an angle of 𝜃 with the positive x axis.},
  width=5.67in,height=3.3in,keepaspectratio]{fig060105}
\color{blue}
\caption{$R=\protect\sqrt{c_1^2+c_2^2}$;  \quad
$c_1=R\cos\theta$;  \quad $c_2=R\sin\theta$}
  \label{figure:6.1.5}
\end{figure}




The motion described by \eqref{eq:6.1.7} or \eqref{eq:6.1.10}
is  \emph{\color{blue}simple harmonic motion}. We see from either of these
equations that the motion is periodic, with period
\[
T=2\pi/\omega_0.
\]
This is the time required for the object to complete one full cycle of
oscillation (for example, to move from its highest position to its
lowest position and back to its highest position). Since the highest
and lowest positions of the object are $y=R$ and $y=-R$, we say that
$R$ is the \emph{\color{blue}amplitude}\index{amplitude!of oscillation} of the oscillation. The angle $\phi$ in
\eqref{eq:6.1.10} is the \emph{\color{blue}phase angle}. It's measured
in
radians. \zcref[S]{eq:6.1.10} is the \indexterm{amplitude-phase
form} of the displacement. If $t$ is in seconds then $\omega_0$ is
in radians per second (rad/s);   it's  the \emph{\color{blue}frequency} of
the motion. It is also called the \emph{\color{blue}natural frequency} of
the spring-mass system without damping.

\begin{example}\label{example:6.1.2}
We found the displacement of the object
 in \zcref{example:6.1.1} to be
\[
y=\frac{3}{2}\cos8t-\frac{3}{8}\sin8t.
\]
Find the frequency, period, amplitude, and phase angle of the motion.
\end{example}

\begin{solution}
The frequency is $\omega_0=\qty{8}{\radian\per\second}$, and the period is
$T=2\pi/\omega_0=\pi/4$ \unit{\second}. Since $c_1=3/2$ and $c_2=-3/8$,
the amplitude is
\[
R=\sqrt{c^2_1+c^2_2}=\sqrt{\left(\frac{3}{2}\right)^2+\left(\frac{3}{8}\right)^2}
=\frac{3}{8}\sqrt{17}.
\]
The phase angle is determined by
\begin{equation}\label{eq:6.1.11}
\cos\phi=\frac{\frac{3}{2}}{\frac{3}{8}\sqrt{17}}=\frac{4}{\sqrt{17}}
\end{equation}
and
\begin{equation}\label{eq:6.1.12}
\sin\phi=\frac{-\frac{3}{8}}{\frac{3}{8}\sqrt{17}}=-\frac{1}{\sqrt{17}}.
\end{equation}
 Using a calculator, we see from \eqref{eq:6.1.11} that
\[
\phi\approx\pm\qty{.245}{\radian}.
\]
Since $\sin\phi<0$ (see \eqref{eq:6.1.12}), the minus sign applies here;
that is,
\[
\phi\approx\qty{-.245}{\radian}.\qedhere
\]
\end{solution}


\begin{example}\label{example:6.1.3}
The natural length of a spring is \qty{1}{\meter}. An object is attached to it and
the length of the spring increases to \qty{102}{\centi\meter} when the object is in
equilibrium. Then the object is initially displaced downward \qty{1}{\centi\meter} and
given an upward velocity of \qty{14}{\centi\meter\per\second}. Find the displacement for
$t>0$. Also, find the natural frequency, period, amplitude, and phase
angle of the resulting motion. Express the answers in cgs units.
\end{example}

\begin{solution}
In cgs units $g=\qty{980}{\centi\meter\per\second\squared}$. Since $\Delta l=\qty{2}{\centi\meter}$,
$\omega_0^2=g/\Delta l=490$.
Therefore
\[
y''+490y=0, \quad  y(0)=-1,\quad y'(0)=14.
\]
 The general solution of the differential equation is
\[
y=c_1\cos7\sqrt{10}t+c_2\sin7\sqrt{10}t,
\]
so
\[
y'=7\sqrt{10}\left(-c_1\sin7\sqrt{10}t+c_2\cos7\sqrt{10}t\right).
\]
Substituting the initial conditions into the last two equations yields
$c_1=-1$ and $c_2=2/\sqrt{10}$. Hence,
\[
y=-\cos7\sqrt{10}t+\frac{2}{\sqrt{10}}\sin7\sqrt{10}t.
\]
The frequency is $7\sqrt{10}$ rad/s, and  the period
is $T=2\pi/(7\sqrt{10})$  s.
The amplitude  is
\[
R=\sqrt{c_1^2+c_2^2}=\sqrt{(-1)^2+\left(\frac{2}{\sqrt{10}}\right)^2}=
\sqrt{\frac{7}{5}}\unit{\centi\meter}.
\]
The  phase angle  is determined by
\[
\cos\phi=\frac{c_1}{R}=-\sqrt{\frac{5}{7}}\quad\text{and}\quad \sin\phi=
\frac{c_2}{R}=\sqrt{\frac{2}{7}}.
\]
Therefore $\phi$ is in the second quadrant and
\[
\phi=\cos^{-1}\left(-\sqrt{\frac{5}{7}}\right)\approx\qty{2.58}{\radian}.\qedhere
\]
\end{solution}


\boxit{Undamped Forced Oscillation}
\index{undamped oscillation!forced}
\index{forced oscillation!undamped}
\index{oscillation!undamped forced}

In many mechanical problems a device is subjected to periodic external
forces. For example, soldiers marching in cadence on a bridge cause
periodic disturbances in the bridge, and the engines of a propeller
driven aircraft cause periodic disturbances in its wings. In the
absence of sufficient damping forces, such disturbances --- even if
small in magnitude --- can cause structural breakdown if they are at
certain critical frequencies. To illustrate, this we'll consider the
motion of an object in a spring-mass system without damping, subject
to an external force
\[
F(t)=F_0\cos\omega t
\]
where $F_0$ is a constant. In this case the equation of motion
\eqref{eq:6.1.2} is
\[
my''+ky=F_0\cos\omega t,
\]
which we rewrite  as
\begin{equation}\label{eq:6.1.13}
y''+\omega_0^2y=\frac{F_0}{m}\cos\omega t
\end{equation}
with $\omega_0=\sqrt{k/m}$.
We'll see from the next two examples that the solutions of
\eqref{eq:6.1.13}
with $\omega\ne\omega_0$ behave very differently from the solutions with
$\omega=\omega_0$.




\begin{example}\label{example:6.1.4}
 Solve the initial value problem
\begin{equation}\label{eq:6.1.14}
y''+\omega_0^2y=\frac{F_0}{m}\cos\omega t, \quad  y(0)=0,\quad y'(0)=0,
\end{equation}
 given that $\omega\ne\omega_0$.
\end{example}


\begin{solution}
We first obtain a particular solution of \eqref{eq:6.1.13} by the method
of undetermined coefficients. Since $\omega\ne\omega _0$,
 $\cos\omega t$ isn't  a solution of the complementary equation
\[
y''+\omega_0^2y=0.
\]
Therefore \eqref{eq:6.1.13} has a particular solution of the form
\[
y_p=A\cos\omega t+B\sin\omega t.
\]
Since
\[
y''_p=-\omega^2(A\cos\omega t+B\sin\omega t),
\]
\[
y_p''+\omega_0^2y_p=\frac{F_0}{m}\cos\omega t
\]
if and only if
\[
(\omega_0^2-\omega^2)\left(A\cos\omega t+
B\sin\omega t\right)=\frac{F_0}{m}\cos\omega t.
\]
 This holds if and only if
\[
A=\frac{F_0}{m(\omega_0^2-\omega^2)}\quad\text{and}\quad  B=0,
\]
so
\[
y_p=\frac{F_0}{m(\omega_0^2-\omega^2)}\cos\omega t.
\]
The general solution of \eqref{eq:6.1.13}  is
\begin{equation}\label{eq:6.1.15}
y=\frac{F_0}{m(\omega_0^2-\omega^2)}\cos\omega t
+c_1\cos\omega_0 t+c_2\sin\omega_0t,
\end{equation}
so
\[
y'=\frac{-\omega F_0}{m(\omega_0^2-\omega^2)}\sin\omega t
+\omega_0(-c_1\sin\omega_0 t+c_2\cos\omega_0t).
\]
The initial conditions $y(0)=0$ and $y'(0)=0$ in \eqref{eq:6.1.14} imply
that
\[
c_1=-\frac{F_0}{m(\omega_0^2-\omega^2)}\quad\text{and}\quad 
c_2=0.
\]
Substituting these into \eqref{eq:6.1.15} yields
\begin{equation}\label{eq:6.1.16}
y=\frac{F_0}{m(\omega_0^2-\omega^2)}(\cos\omega t-\cos\omega_0t).\qedhere
\end{equation}
\end{solution}

It is revealing to write this in a different form.  We start with the
trigonometric identities
\begin{align*}
\cos(\alpha-\beta)&=\cos\alpha\cos\beta+\sin\alpha\sin\beta\\
\cos(\alpha+\beta)&=\cos\alpha\cos\beta-\sin\alpha\sin\beta.
\end{align*}
Subtracting the second identity from the first yields
\begin{equation}\label{eq:6.1.17}
\cos(\alpha-\beta)-\cos(\alpha+\beta)=2\sin\alpha\sin\beta
\end{equation}
 Now let
\begin{equation}\label{eq:6.1.18}
\alpha-\beta=\omega t\quad\text{and}\quad \alpha+\beta=\omega_0t,
\end{equation}
 so that
\begin{equation}\label{eq:6.1.19}
\alpha=\frac{(\omega_0+\omega)t}{2}\quad\text{and}\quad
\beta=\frac{(\omega_0-\omega)t}{2}.
\end{equation}
 Substituting \eqref{eq:6.1.18} and \eqref{eq:6.1.19} into \eqref{eq:6.1.17}
yields
 \[
\cos\omega t-\cos\omega_0t=
2\sin\frac{(\omega_0-\omega)t}{2}\sin\frac{(\omega_0+\omega)t}{2},
\]
and substituting this into \eqref{eq:6.1.16} yields
\begin{equation}\label{eq:6.1.20}
y=R(t)\sin\frac{(\omega_0+\omega)t}{2},
\end{equation}
 where
\begin{equation}\label{eq:6.1.21}
R(t)=\frac{2F_0}{m(\omega_0^2-\omega^2)}
\sin\frac{(\omega_0-\omega)t}{2}.
\end{equation}

From \eqref{eq:6.1.20} we can regard $y$ as a sinusoidal variation with
frequency $(\omega_0+\omega)/2$ and variable amplitude
$|R(t)|$. In \zcref{figure:6.1.6} the dashed curve above the $t$ axis
is $y=|R(t)|$, the dashed curve below the $t$ axis is $y=-|R(t)|$, and
the displacement $y$ appears as an oscillation bounded by them. The
oscillation of $y$ for $t$ on an interval between successive zeros of
$R(t)$ is called a \indexterm{beat}.



\begin{figure}
  \centering
  \includegraphics[alt={A high frequency sine-like curve whose amplitude is given by a lower frequency sinusoidal.},
  width=5.67in,height=3.3in,keepaspectratio]{fig060106}
\color{blue}
\caption{Undamped oscillation with beats}
  \label{figure:6.1.6}
\end{figure}




You can see from
\eqref{eq:6.1.20} and \eqref{eq:6.1.21} that
\[
|y(t)|\le\frac{2|F_0|}{m|\omega_0^2-\omega^2|};
\]
moreover, if $\omega+\omega_0$ is sufficiently large compared with $\omega
-\omega_0$, then $|y|$ assumes values close to (perhaps equal to) this
upper bound during each beat. However, the oscillation remains bounded for
all $t$. (This assumes that the spring can withstand
deflections of this size and continue to obey Hooke's law.)  The next
example shows that this isn't  so if $\omega=\omega_0$.

\begin{example}\label{example:6.1.5}
Find the general solution of
\begin{equation}\label{eq:6.1.22}
y''+\omega_0^2y=\frac{F_0}{m}\cos\omega_0t.
\end{equation}
\end{example}


\begin{solution}
We first obtain a particular solution $y_p$ of \eqref{eq:6.1.22}. Since
$\cos\omega_0t$ is a solution of the complementary equation, the form
for $y_p$ is
\begin{equation}\label{eq:6.1.23}
y_p=t(A\cos\omega_0t+B\sin\omega_0t).
\end{equation}
 Then
\[
y_p'=A\cos\omega_0t+B\sin\omega_0t
+\omega_0t(-A\sin\omega_0t+B\cos\omega_0t)
\]
and
\begin{equation}\label{eq:6.1.24}
y''_p=2\omega_0(-A\sin\omega_0t
+B\cos\omega_0t)-\omega_0^2t(A\cos\omega_0t+B\sin\omega_0t).
\end{equation}
From \eqref{eq:6.1.23} and \eqref{eq:6.1.24}, we see that $y_p$ satisfies
\eqref{eq:6.1.22} if
\[
-2A\omega_0\sin\omega_0t+2B\omega_0\cos\omega_0t=\frac{F_0}{m}
\cos\omega_0t;
\]
that is, if
\[
A=0\quad\text{and}\quad  B=\frac{F_0}{2m\omega_0}.
\]
Therefore
\[
y_p=\frac{F_0t}{2m\omega_0}\sin\omega_0t
\]
is a particular solution of \eqref{eq:6.1.22}. The general solution of
 \eqref{eq:6.1.22} is
\[
y=\frac{F_0t}{2m\omega_0}\sin\omega_0t+c_1\cos\omega_0t+c_2\sin\omega_0t.
\]
The graph of $y_p$ is shown in \zcref{figure:6.1.7},  where it can be
seen that $y_p$ oscillates between the dashed lines
\[
y=\frac{F_0t}{2m\omega_0}\quad\text{and}\quad  y=-\frac{F_0t}{2m\omega_0}
\]
with increasing amplitude that approaches $\infty$ as $t\to\infty$. Of
course, this means that the spring must eventually fail to obey Hooke's law
or break.
\end{solution}


This phenomenon of unbounded displacements of a spring-mass system in
response to a periodic forcing function at its natural
frequency is called \emph{\color{blue}resonance}. More complicated mechanical
structures can also exhibit resonance-like phenomena. For example,
rhythmic oscillations of a suspension bridge by wind forces or of an
airplane wing by periodic vibrations of reciprocating engines can
cause damage or even failure if the frequencies of the disturbances
are close to critical frequencies determined by the parameters of the
mechanical system in question.



\begin{figure}
  \centering
  \includegraphics[alt={A sine-like curve whose amplitude grows linearly.},
  width=5.67in,height=3.3in,keepaspectratio]{fig060107}
\color{blue}
\caption{Unbounded displacement due to resonance}
  \label{figure:6.1.7}
\end{figure}



\exercises
 In the following exercises assume that there's no damping.

\begin{exerciselist}

\item\label{exer:6.1.1}\CGex
An object stretches a spring 4 inches in equilibrium.
Find and graph its displacement for $t>0$ if it's initially displaced
36 inches above equilibrium and given a downward velocity of \qty{2}{\feet\per\second}.

\item\label{exer:6.1.2}An object  stretches a string 1.2 inches in equilibrium.
Find its displacement for $t>0$ if it's initially displaced 3 inches
below equilibrium and given a downward velocity of \qty{2}{\feet\per\second}.

\item\label{exer:6.1.3}A spring with natural length \qty{.5}{\meter} has length \qty{50.5}{\centi\meter} with a mass of \qty{2}{\gram} suspended from it. The mass is initially displaced \qty{1.5}{\centi\meter} below
equilibrium and released with zero velocity. Find its displacement for
$t>0$.

\item\label{exer:6.1.4}An object stretches a spring 6 inches in equilibrium. Find its
displacement for $t>0$ if it's initially displaced 3 inches above
equilibrium and given a downward velocity of 6 inches/s. Find the
frequency, period, amplitude and phase angle of the motion.

\item\label{exer:6.1.5}\CGex
An object stretches a spring \qty{5}{\centi\meter} in equilibrium. It is initially
displaced \qty{10}{\centi\meter} above equilibrium and given an upward velocity of \qty{.25}{\meter\per\second}. Find and graph its displacement for $t>0$. Find the frequency,
period, amplitude, and phase angle of the motion.

\item\label{exer:6.1.6}A \qty{10}{\kilo\gram} mass stretches a spring \qty{70}{\centi\meter} in equilibrium. Suppose a \qty{2}{\kilo\gram} mass is
attached to the spring, initially displaced \qty{25}{\centi\meter} below equilibrium,
and given an upward velocity of \qty{2}{\meter\per\second}. Find its displacement for
$t>0$. Find the frequency, period, amplitude, and phase angle of the
motion.

\item\label{exer:6.1.7}A  weight stretches a spring 1.5 inches in equilibrium. The weight
is initially displaced 8 inches above equilibrium and given a downward
velocity of \qty{4}{\feet\per\second}. Find its displacement for $t > 0$.

\item\label{exer:6.1.8}A weight stretches a spring 6 inches in equilibrium. The weight is
initially displaced 6 inches above equilibrium and given a downward
velocity of \qty{3}{\feet\per\second}. Find its displacement for $t>0$.

\item\label{exer:6.1.9}A spring-mass system has natural frequency $7\sqrt{10}$ \unit{\radian\per\second}. The
natural length of the spring is \qty{.7}{\meter}. What is the length of the spring
when the mass is in equilibrium?

\item\label{exer:6.1.10}A \qty{64}{\pound} weight is attached to a spring with constant $k=\qty{8}{\pound\per\foot}$ and
subjected to an external force $F(t)=2\sin t$. The weight is initially
displaced 6 inches above equilibrium and given an upward velocity of \qty{2}{\feet\per\second}. Find its displacement for $t>0$.

\item\label{exer:6.1.11}A unit mass hangs in equilibrium from a spring with constant $k=1/16$.
Starting at $t=0$, a force $F(t)=3\sin t$ is applied to the mass. Find
its displacement for $t>0$.

\item\label{exer:6.1.12}\CGex
A \qty{4}{\pound} weight stretches a spring \qty{1}{\foot} in equilibrium. An external
force $F(t)=.25\sin8 t$ \unit{\pound} is applied to the weight, which is
initially displaced 4 inches above equilibrium and given a downward
velocity of \qty{1}{\foot\per\second}. Find and graph its displacement for $t>0$.

\item\label{exer:6.1.13}A \qty{2}{\pound} weight stretches a spring 6 inches in equilibrium.  An external
force $F(t)=\sin8t$ \unit{\pound} is applied to the weight, which is released
from rest 2 inches below equilibrium. Find its displacement
for $t>0$.

\item\label{exer:6.1.14}A \qty{10}{\gram} mass suspended on a spring moves in simple harmonic
motion with period \qty{4}{\second}.  Find the period of the simple
harmonic motion of a \qty{20}{\gram} mass suspended from the same spring.


\item\label{exer:6.1.15}A \qty{6}{\pound} weight stretches a spring 6 inches in equilibrium. Suppose
an external force $F(t)=\dfrac{3}{16}\sin\omega
t+\dfrac{3}{8}\cos\omega t $ \unit{\pound} is applied to the weight. For what
value
of $\omega$ will the displacement be unbounded? Find the displacement
if $\omega$ has this value. Assume that the motion starts from
equilibrium with zero initial velocity.

\item\label{exer:6.1.16}\CGex
A \qty{6}{\pound} weight stretches a spring 4 inches in equilibrium. Suppose
an external force $ F(t)=4\sin\omega t-6\cos\omega t $ \unit{\pound} is applied
to the weight. For what value of $\omega$ will the displacement be
unbounded? Find and graph the displacement if $\omega$ has this value.
Assume that the motion starts from equilibrium with zero initial
velocity.

\item\label{exer:6.1.17}A mass of one \unit{\kilo\gram} is attached to a spring with constant $k=\qty{4}{\newton\per\meter}$. An
external force $F(t)=-\cos\omega t-2\sin\omega t$ \unit{\newton} is applied to the
mass. Find the displacement $y$ for $t>0$ if $\omega$ equals the
natural frequency of the spring-mass system. Assume that the
mass is initially displaced \qty{3}{\meter} above equilibrium and given an upward
velocity of \qty{450}{\centi\meter\per\second}.

\item\label{exer:6.1.18}An object is in simple harmonic motion
with frequency $\omega_0$, with $y(0)=y_0$ and $y'(0)=v_0$. Find
its displacement for $t>0$. Also, find the amplitude of the
oscillation and give formulas for the sine and cosine of the initial
phase angle.

\item\label{exer:6.1.19}Two objects suspended from identical springs are set into
motion.  The period of one object is twice the period of the other.
How are the weights of the two objects related?

\item\label{exer:6.1.20}Two objects suspended from identical springs are set into motion. The
weight of one object is twice the weight of the other. How are the
periods of the resulting motions related?

\item\label{exer:6.1.21}Two identical objects suspended from different springs are set into
motion. The period of one motion is 3 times the period of the
other. How are the two spring constants related?
\end{exerciselist}



\section{Spring Problems II}\label{sec:spring_ii}
\index{spring problems}

\boxit{Free Vibrations With Damping}
\index{free vibrations with damping}
\index{damping!free vibrations with}
\index{vibrations!free with damping}

In this section we consider the motion of an object in a spring-mass
system with damping. We start with unforced motion, so the equation of
motion is
\begin{equation}\label{eq:6.2.1}
my''+cy'+ky=0.
\end{equation}
Now suppose   the object is displaced from equilibrium and given an
initial velocity. Intuition suggests that if the damping force is
sufficiently weak the resulting motion will be oscillatory, as in the
undamped case considered in the previous section, while if it's
sufficiently strong the object may just move slowly toward the
equilibrium
position without ever reaching it. We'll now confirm these intuitive
ideas mathematically. The characteristic equation of \eqref{eq:6.2.1} is
\[
mr^2+cr+k=0.
\]
The roots of this equation are
\begin{equation}\label{eq:6.2.2}
r_1=\frac{-c-\sqrt{c^2-4mk}}{2m}\quad\text{and}\quad  r_2=
\frac{-c+\sqrt{c^2-4mk}}{2m}.
\end{equation}
In \zcref{sec:const_coeff_homog} we saw that the form of the
solution of  \eqref{eq:6.2.1} depends upon whether $c^2-4mk$ is positive,
negative, or zero. We'll now consider these three cases.


\boxit{Underdamped Motion}
\index{underdamped motion}
\index{motion!underdamped}

We say the motion is \emph{\color{blue}underdamped} if $c<\sqrt{4mk}$. In this
case $r_1$ and $r_2$ in \eqref{eq:6.2.2} are complex conjugates, which we
write as
\[
r_1=-\frac{c}{2m}-i\omega_1\quad\text{and}\quad 
r_2=-\frac{c}{2m}+i\omega_1,
\]
where
\[
\omega_1=\frac{\sqrt{4mk-c^2}}{2m}.
\]
The general solution of  \eqref{eq:6.2.1} in this case is
\[
y=e^{-ct/2m}(c_1\cos\omega_1 t+c_2\sin\omega_1 t).
\]
By the method used in \zcref{sec:spring_i} to derive the
amplitude-phase form of the displacement of an object in simple
harmonic motion, we can rewrite this equation as
\begin{equation}\label{eq:6.2.3}
y=Re^{-ct/2m}\cos(\omega_1 t-\phi),
\end{equation}
 where
\[
R=\sqrt{c_1^2+c_2^2},\quad R\cos\phi=c_1,\quad\text{and}\quad  R\sin\phi=c_2.
\]
The factor $Re^{-ct/2m}$ in \eqref{eq:6.2.3} is called the
\indexterm{time-varying amplitude}\index{amplitude!time-varying} of the motion, the quantity $\omega_1$ is
called the \indexterm{frequency}, and $T=2\pi/\omega_1$ (which is the
period of the cosine function in \eqref{eq:6.2.3} is called the
\indexterm{quasi-period}. A typical graph of \eqref{eq:6.2.3} is shown in
\zcref{figure:6.2.1}. As illustrated in that figure, the graph of $y$
oscillates between the dashed exponential curves $y=\pm Re^{-ct/2m}$.

\begin{figure}
  \centering
  \includegraphics[alt={A sine-like curve whose amplitude is exponentially decreasing.},
  width=5.67in,height=3.3in,keepaspectratio]{fig060201}
\color{blue}
\caption{ Underdamped motion}
  \label{figure:6.2.1}
\end{figure}


\boxit{Overdamped Motion}
\index{overdamped motion}
\index{motion!overdamped}

We say the motion is \emph{\color{blue}overdamped} if $c>\sqrt{4mk}$. In this
case the zeros $r_1$ and $r_2$ of the characteristic polynomial are
real, with $r_1<r_2<0$ (see \eqref{eq:6.2.2}), and the general solution of
\eqref{eq:6.2.1} is
\[
y=c_1e^{r_1t}+c_2e^{r_2t}.
\]
Again $\lim_{t\to\infty}y(t)=0$ as in the
underdamped case, but the motion isn't  oscillatory, since $y$ can't
equal zero for more than one value of $t$ unless $c_1=c_2=0$.
(\zcref{exer:6.2.23}.)


\boxit{Critically Damped Motion}
\index{critically damped motion}
\index{motion!critically damped}

We say the motion is \emph{\color{blue}critically damped} if $c=\sqrt{4mk}$. In
this case $r_1=r_2=-c/2m$ and the general solution of \eqref{eq:6.2.1} is
\[
y=e^{-ct/2m}(c_1+c_2t).
\]

Again $\lim_{t\to\infty}y(t)=0$ and the motion is nonoscillatory,
since $y$ can't equal zero for more than one value of $t$ unless
$c_1=c_2=0$ (\zcref{exer:6.2.22}).

\begin{example}\label{example:6.2.1}
Suppose a \qty{64}{\pound} weight stretches a spring 6 inches in equilibrium
and a dashpot provides a damping force of $c$ \unit{\pound} for each \unit{\foot\per\second} of
velocity.
\begin{alist}
\item % (a)
 Write the equation of motion of the object and determine the value of
$c$ for which the motion is critically damped.
\item % (b)
 Find the displacement $y$ for $t>0$ if the motion is critically damped
and the initial conditions are $y(0)=1$ and $y'(0)=20$.
\item % (c)
 Find the displacement $y$ for $t>0$ if the motion is critically damped
and the initial conditions are $y(0)=1$ and $y'(0)=-20$.
\end{alist}
\end{example}

\begin{alistSolution}
\item Here $m=2$ slugs and $k=64/.5=\qty{128}{\pound\per\feet}$.
Therefore the equation of motion  \eqref{eq:6.2.1} is
\begin{equation}\label{eq:6.2.4}
2y''+cy'+128y=0.
\end{equation}
The characteristic equation  is
\[
2r^2+cr+128=0,
\]
which has  roots
\[
r=\frac{-c\pm\sqrt{c^2-8\cdot128}}{4}.
\]
Therefore the damping is critical if
\[
c=\sqrt{8\cdot128}=\qty{32}{\pound\second\per\foot}.
\]

\item
 Setting $c=32$ in  \eqref{eq:6.2.4} and cancelling the common factor $2$
yields
\[
y''+16y+64y=0.
\]
The characteristic equation  is
\[
r^2+16r+64y=(r+8)^2=0.
\]
Hence, the general solution is
\begin{equation}\label{eq:6.2.5}
y=e^{-8t}(c_1+c_2t).
\end{equation}
 Differentiating this yields
\begin{equation}\label{eq:6.2.6}
y'=-8y+c_2e^{-8t}.
\end{equation}
Imposing the initial conditions $y(0)=1$ and $y'(0)=20$ in the last
two equations shows that $1=c_1$ and $20=-8+c_2$. Hence, the solution
of the initial value problem is
\[
y=e^{-8t}(1+28t).
\]
Therefore the object approaches equilibrium from above as
$t\to\infty$. There's no oscillation.

\item
Imposing the initial conditions $y(0)=1$ and $y'(0)=-20$ in
\eqref{eq:6.2.5} and \eqref{eq:6.2.6} yields $1=c_1$ and $-20=-8+c_2$. Hence,
the solution of this initial value problem is
\[
y=e^{-8t}(1-12t).
\]
Therefore the object moves downward through equilibrium just once,
and then approaches equilibrium from below as $t\to\infty$. Again,
there's no oscillation. The solutions of these two initial value
problems are graphed in \zcref{figure:6.2.2}.

\begin{figure}
  \centering
  \includegraphics[alt={Two curves that begin at (0,1).  Curve (a) increases and then decreases exponentially.  Curve (b) decreases to cross the x axis near 0.1 and then decays exponentially to 0.},
  width=5.67in,height=3.3in,keepaspectratio]{fig060202}
\color{blue}
\caption{(a) $y=e^{-8t}(1+28t)$, (b) $y=e^{-8t}(1-12t)$}
  \label{figure:6.2.2}
\end{figure}

\end{alistSolution}




\begin{example}
\label{example:6.2.2}
Find the displacement of the object in \zcref{example:6.2.1} if the
damping constant is $c=\qty{4}{\pound\second\per\feet}$ and the initial conditions are
$y(0)=\qty{1.5}{\feet}$ and $y'(0)=\qty{-3}{\feet\per\second}$.
\end{example}


\begin{solution}
With $c=4$, the equation of motion \eqref{eq:6.2.4} becomes
\begin{equation}\label{eq:6.2.7}
y''+2y'+64y=0
\end{equation}
after cancelling the common factor 2. The characteristic equation
\[
r^2+2r+64=0
\]
has complex conjugate roots
\[
r=\frac{-2\pm\sqrt{4-4\cdot64}}{2}=-1\pm3\sqrt7i.
\]
Therefore the motion is underdamped and the
general solution of  \eqref{eq:6.2.7} is
\[
y=e^{-t}(c_1\cos3\sqrt7t+c_2\sin3\sqrt7t).
\]
Differentiating this yields
\[
y'=-y+3\sqrt7e^{-t}(-c_1\sin3\sqrt7t+c_2\cos3\sqrt7t).
\]
Imposing the initial conditions $y(0)=1.5$ and $y'(0)=-3$ in the last
two equations yields $1.5=c_1$ and $-3=-1.5+3\sqrt7c_2$. Hence, the
solution of the initial value problem is
\begin{equation}\label{eq:6.2.8}
y=e^{-t}\bigg(\frac{3}{2}\cos3\sqrt7t-\frac{1}{2\sqrt7}
\sin3\sqrt7t\bigg).
\end{equation}
 The amplitude of the function in parentheses is
\[
 R=\sqrt{\bigg(\frac{3}{2}\bigg)^2+\bigg(\frac{1}{2\sqrt7}\bigg)^2}
=\sqrt{\frac{9}{4}+\frac{1}{4\cdot7}}
=\sqrt\frac{64}{4\cdot7}=\frac{4}{\sqrt7}.
\]
Therefore we can rewrite  \eqref{eq:6.2.8} as
\[
y=\frac{4}{\sqrt7}e^{-t}\cos(3\sqrt7t-\phi),
\]
 where
\[
\cos\phi=\frac{3}{2R}=\frac{3\sqrt7}{8}\quad\text{and}\quad \sin\phi=-\frac{1}{2\sqrt7R}=
-\frac{1}{8}.
\]
Therefore $\phi\cong\qty{-.125}{\radian}$.
\end{solution}

\begin{example}\label{example:6.2.3}
Let the damping constant in \zcref{example:6.2.1} be $c=\qty{40}{\pound\second\per\foot}$. Find the
displacement $y$ for $t>0$ if $y(0)=1$ and $y'(0)=1$.
\end{example}

\begin{solution}
With $c=40$, the equation of motion \eqref{eq:6.2.4} reduces to
\begin{equation}\label{eq:6.2.9}
y''+20y'+64y=0
\end{equation}
 after cancelling the  common factor 2.
The characteristic equation
\[
r^2+20r+64=(r+16)(r+4)=0
\]
 has the roots $r_1=-4$ and $r_2=-16$.  Therefore the general
solution of  \eqref{eq:6.2.9} is
\begin{equation}\label{eq:6.2.10}
y=c_1e^{-4t}+c_2e^{-16t}.
\end{equation}
 Differentiating this yields
\[
y'=-4e^{-4t}-16c_2e^{-16t}.
\]
The last two equations and the initial conditions $y(0)=1$ and $y'(0)=1$
imply that
\begin{align*}
c_1+\phantom{16}c_2&=1\\
-4c_1-16c_2&=1.
\end{align*}
 The solution of this system is $c_1=17/12$, $c_2=-5/12$.
Substituting these into  \eqref{eq:6.2.10} yields
\[
y=\frac{17}{12}e^{-4t}-\frac{5}{12}e^{-16t}
\]
as the solution of the given initial value problem
(\zcref{figure:6.2.3}).
\end{solution}

\begin{figure}
  \centering
  \includegraphics[alt={A curve beginning at (0,1), increasing for just a bit, and then decreasing exponentially toward 0.},
  width=5.67in,height=3.3in,keepaspectratio]{fig060203}
\color{blue}
\caption{$y=\dfrac{17}{12}e^{-4t}-\dfrac{5}{12}e^{-16t}$}
  \label{figure:6.2.3}
\end{figure}


\boxit{Forced Vibrations With Damping}
\index{forced vibrations with damping}
\index{damping!forced vibrations with}
\index{vibrations!forced with damping}

Now we consider the motion of an object in a spring-mass system with
damping, under the influence of a periodic forcing function
$F(t)=F_0\cos\omega t$, so that the equation of motion is
\begin{equation}\label{eq:6.2.11}
my''+cy'+ky=F_0\cos\omega t.
\end{equation}
In \zcref{sec:spring_i} we considered this equation with $c=0$ and
found that the resulting displacement $y$ assumed arbitrarily large
values in the case of resonance   (that is, when
$\omega=\omega_0=\sqrt{k/m}$). Here we'll see that in the presence
of damping the displacement remains bounded for all $t$, and the
initial conditions have little effect on the motion as $t\to\infty$.
In fact, we'll see that for large $t$ the displacement is closely
approximated by a function of the form
\begin{equation}\label{eq:6.2.12}
y=R\cos(\omega t-\phi),
\end{equation}
where the amplitude $R$ depends upon $m$, $c$, $k$, $F_0$, and
$\omega$. We're interested in the following question:

\bigbox{\textsc{Question:}\slshape Assuming that $m$, $c$, $k$, and $F_0$ are
held constant, what value of $\omega$ produces the largest amplitude
$R$ in \eqref{eq:6.2.12}, and what is this largest amplitude?} To
answer this question, we must solve \eqref{eq:6.2.11} and determine $R$ in
terms of $F_0,\omega_0,\omega$, and $c$. We can obtain a particular
solution of \eqref{eq:6.2.11} by the method of undetermined coefficients.
Since $\cos\omega t$ does not satisfy the complementary equation
\[
my''+cy'+ky=0,
\]
we can obtain a particular solution of \eqref{eq:6.2.11} in the form
\begin{equation}\label{eq:6.2.13}
y_p=A\cos\omega t+B\sin\omega t.
\end{equation}
Differentiating this yields
\[
y_p'=\quad\omega (-A\sin\omega t+B\cos\omega t)
\]
and
\[
y_p''=-\omega^2(A\cos\omega t+B\sin\omega t).
\]
 From the last three equations,
\[
my''_p+cy'_p+ky_p=(-m\omega^2A+c\omega B+kA)\cos\omega t+
(-m\omega^2 B-c\omega A+kB)\sin\omega t,
\]
so $y_p$ satisfies  \eqref{eq:6.2.11}  if
\begin{alignat*}{2}
(k-m\omega^2) A&+ c\omega B &&=F_0\\
-c\omega A&+(k-m\omega^2)B &&=0.
\end{alignat*}

Solving  for $A$ and $B$
and substituting the results into  \eqref{eq:6.2.13} yields
\[
y_p=\frac{F_0}{(k-m\omega^2)^2+c^2\omega^2}
\left[(k-m\omega^2)\cos\omega t+c\omega\sin\omega t\right],
\]
which can be written in amplitude-phase form  as
\begin{equation}\label{eq:6.2.14} %\springsteady
y_p=\frac{F_0}{\sqrt{(k-m\omega^2)^2+c^2\omega^2}}
\cos(\omega t-\phi),
\end{equation}
where
\begin{equation}\label{eq:6.2.15}
\cos\phi=\frac{k-m\omega^2}{\sqrt
{(k-m\omega^2)^2+c^2\omega^2}}\quad\text{and}\quad 
\sin\phi=\frac{c\omega}{\sqrt{(k-m\omega^2)^2+c^2\omega^2}}.
\end{equation}

  To compare this with the undamped forced vibration
that we considered in \zcref{sec:spring_i} it's useful to write
\begin{equation}\label{eq:6.2.16}
k-m\omega^2=m\bigg(\frac{k}{m}-\omega^2\bigg)=
m(\omega_0^2-\omega^2),
\end{equation}
 where $\omega_0=\sqrt{k/m}$ is the natural
angular frequency of the undamped simple harmonic motion of an
object with mass $m$ on a spring with constant $k$.
Substituting  \eqref{eq:6.2.16} into  \eqref{eq:6.2.14} yields
\begin{equation}\label{eq:6.2.17}
y_p=\frac{F_0}{\sqrt{m^2(\omega^2_0-\omega^2)^2+
c^2\omega^2}}\cos(\omega t-\phi).
\end{equation}

The solution of an initial value problem
\[
my''+cy'+ky=F_0\cos\omega t, \quad  y(0)=y_0,\quad y'(0)=v_0,
\]
is of the form $y=y_c+y_p$,
where  $y_c$ has one of the three forms
\begin{align*}
y_c&=e^{-ct/2m}(c_1\cos\omega_1t+c_2\sin\omega_1t),\\
y_c&=e^{-ct/2m}(c_1+c_2t),\\
y_c&=c_1e^{r_1t}+c_2e^{r_2t}\,(r_1,r_2<0).
\end{align*}
In all three cases $\lim_{t\to\infty} y_c(t)=0$ for any choice of
$c_1$ and $c_2$. For this reason we say that $y_c$ is the
\emph{\color{blue}transient component}\index{transient!component} of the solution $y$.
The behavior of $y$ for
large $t$ is determined by $y_p$, which we call the \emph{\color{blue}steady state
component}\index{steady state!component} of $y$. Thus, for large $t$ the motion is like simple
harmonic motion at the frequency of the external force.

The amplitude of $y_p$ in \eqref{eq:6.2.17} is
\begin{equation}\label{eq:6.2.18}
R=\frac{F_0}{\sqrt{m^2(\omega^2_0-\omega^2)^2+c^2\omega^2}},
\end{equation}
which is finite for all $\omega$; that is, the presence of damping
precludes the phenomenon of resonance that we encountered in studying
undamped vibrations under a periodic forcing function. We'll now
find the value $\omega_{\max}$ of $\omega$ for which $R$ is maximized.
This is the value of $\omega$ for which the function
\[
\rho (\omega)=m^2(\omega^2_0-\omega^2)^2+c^2\omega^2
\]
in the denominator of  \eqref{eq:6.2.18} attains its minimum value.  By
rewriting this as
\begin{equation}\label{eq:6.2.19}
\rho (\omega)=m^2(\omega^4_0+\omega^4)+
 (c^2-2m^2\omega^2_0)\omega^2,
\end{equation}
 you can see that $\rho$ is a strictly
increasing function of $\omega^2$ if
\[
c\ge\sqrt{2m^2\omega^2_0}=\sqrt{2mk}.
\]
(Recall that $\omega^2_0=k/m$). Therefore $\omega_{\max}=0$ if this
inequality holds. From \eqref{eq:6.2.15}, you can see that $\phi=0$ if
$\omega=0$. In this case, \eqref{eq:6.2.14} reduces to
\[
y_p=\frac{F_0}{\sqrt{m^2\omega^4_0}}=\frac{F_0}{k},
\]
which is consistent with Hooke's law: if the mass is subjected to a
constant force $F_0$,  its displacement should approach a constant
$y_p$ such that $ky_p=F_0$. Now suppose   $c<\sqrt{2mk}$. Then,
from \eqref{eq:6.2.19},
\[
\rho'(\omega)=2\omega(2m^2\omega^2+c^2-2m^2\omega^2_0),
\]
and $\omega_{\max}$ is the value of $\omega$ for which the
expression in parentheses equals zero; that is,
\[
\omega_{\max}=\sqrt{\omega^2_0-\frac{c^2}{2m^2}}
=\sqrt{\frac{k}{m}\left(1-\frac{c^2}{2km}\right)}.
\]
(To see that $\rho(\omega_{\max})$ is the minimum value of
$\rho(\omega)$, note that $\rho'(\omega)<0$ if $\omega
<\omega_{\max}$ and $\rho'(\omega)>0$ if $\omega>\omega_{\max}$.)
Substituting $\omega=\omega_{\max}$ in \eqref{eq:6.2.18} and simplifying
shows that the maximum amplitude $R_{\max}$ is
\[
R_{\max}=\frac{2mF_0}{c\sqrt{4mk-c^2}} \quad\text{if}\quad c<\sqrt{2mk}.
\]
We summarize our results  as follows.


\begin{theorem}\color{blue} \label{thmtype:6.2.1}
 Suppose we consider the amplitude $R$ of the steady state
component of the solution of
\[
my''+cy'+ky=F_0\cos\omega t
\]
as a function of $\omega$.
\begin{alist}
\item If $c\ge\sqrt{2mk}$,  the maximum amplitude is
 $R_{\max}=F_0/k$ and it's attained when $\omega=
\omega_{\max}=0$.
\item If $c<\sqrt{2mk}$, the maximum amplitude is
\begin{equation}\label{eq:6.2.20}
R_{\max}=\frac{2m F_0}{c\sqrt{4mk-c^2}},
\end{equation}
and it's attained when
\begin{equation}\label{eq:6.2.21}
\omega=\omega_{\max}=\sqrt{\frac{k}{m}\left(1-\frac{c^2}{2km}\right)}.
\end{equation}
\end{alist}
\end{theorem}
Note that $R_{\max}$ and $\omega_{\max}$ are continuous functions
of $c$, for $c\ge0$, since  \eqref{eq:6.2.20} and  \eqref{eq:6.2.21} reduce
to $R_{\max}=F_0/k$ and $\omega_{\max}=0$ if $c=\sqrt{2km}$.

\exercises
\begin{exerciselist}

\item\label{exer:6.2.1}A \qty{64}{\pound} object stretches a spring \qty{4}{\feet} in equilibrium. It is attached
to a dashpot with damping constant $c=\qty{8}{\pound\second\per\foot}$. The object is
initially displaced 18 inches above equilibrium and given a downward
velocity of \qty{4}{\feet\per\second}. Find its displacement and time-varying
amplitude for $t>0$.

\item\label{exer:6.2.2}\CGex
A \qty{16}{\pound} weight is attached to a spring with natural length \qty{5}{\feet}. With
the weight attached, the spring measures \qty{8.2}{\feet}. The weight is
initially displaced \qty{3}{\feet} below equilibrium and given an upward
velocity of \qty{2}{\feet\per\second}. Find and graph its displacement
 for $t>0$ if the medium resists the motion with a force of
one \unit{\pound} for each \unit{\foot\per\second} of velocity. Also, find its time-varying amplitude.

\item\label{exer:6.2.3}\CGex
An \qty{8}{\pound} weight stretches a spring 1.5 inches. It is attached to a
dashpot with damping constant $c=\qty{8}{\pound\second\per\foot}$. The weight is
initially displaced 3 inches above equilibrium and given an upward
velocity of \qty{6}{\feet\per\second}. Find and graph its displacement for $t>0$.

\item\label{exer:6.2.4}A \qty{96}{\pound} weight stretches a spring \qty{3.2}{\feet} in equilibrium. It is
attached to a dashpot with damping constant $c=\qty{18}{\pound\second\per\foot}$. The
weight is initially displaced 15 inches below equilibrium and given a
downward velocity of \qty{12}{\feet\per\second}. Find its displacement for $t>0$.

\item\label{exer:6.2.5}A \qty{16}{\pound} weight  stretches a spring 6 inches in equilibrium.
It is attached to a damping mechanism with constant $c$.  Find all values
of $c$ such that the free vibration of the weight has infinitely many
oscillations.

\item\label{exer:6.2.6}An \qty{8}{\pound} weight stretches a spring \qty{.32}{\feet}. The weight is initially
displaced 6 inches above equilibrium and given an upward velocity of \qty{4}{\feet\per\second}. Find its displacement for $t>0$ if the medium exerts a
damping force of \qty{1.5}{\pound} for each \unit{\foot\per\second} of velocity.

\item\label{exer:6.2.7}A \qty{32}{\pound} weight stretches a spring \qty{2}{\feet} in equilibrium. It is attached
to a dashpot with constant $c=\qty{8}{\pound\second\per\foot}$. The weight is initially
displaced 8 inches below equilibrium and released from rest. Find its
displacement for $t>0$.

\item\label{exer:6.2.8}A mass of \qty{20}{\gram} stretches a  spring \qty{5}{\centi\meter}.  The spring
is attached to a dashpot with damping constant \qty{400}{\dyne\second\per\centi\meter}.  Determine
the displacement for $t>0$ if the mass is  initially displaced
\qty{9}{\centi\meter} above equilibrium and released from rest.

\item\label{exer:6.2.9}A \qty{64}{\pound} weight is suspended from a spring with constant $k=\qty{25}{\pound\per\foot}$.
It is initially displaced 18 inches above equilibrium and released
from rest. Find its displacement for $t>0$ if the medium resists the
motion with \qty{6}{\pound} of force for each \unit{\foot\per\second} of velocity.

\item\label{exer:6.2.10}A \qty{32}{\pound} weight stretches a spring \qty{1}{\foot} in equilibrium. The weight is
initially displaced 6 inches above equilibrium and given a downward
velocity of \qty{3}{\feet\per\second}. Find its displacement for $t>0$ if the medium
resists the motion with a force  equal to 3 times the speed
in \unit{\feet\per\second}.

\item\label{exer:6.2.11}An \qty{8}{\pound} weight stretches a spring 2 inches. It is attached to a
dashpot with damping constant $c=\qty{4}{\pound\second\per\foot}$. The weight is
initially displaced 3 inches above equilibrium and given a downward
velocity of \qty{4}{\feet\per\second}. Find its displacement for $t>0$.

\item\label{exer:6.2.12}\CGex
A \qty{2}{\pound} weight stretches a spring \qty{.32}{\feet}. The weight is initially
displaced 4 inches below equilibrium and given an upward velocity of \qty{5}{\feet\per\second}. The medium provides damping with constant $c=1/8$ \unit{\pound\second\per\foot}.
Find and graph the displacement for $t>0$.

\item\label{exer:6.2.13}An \qty{8}{\pound} weight stretches a spring 8 inches in equilibrium.
It is attached to a dashpot with damping constant $c=\qty{.5}{\pound\second\per\foot}$
and subjected to an external force  $F(t)=4\cos2t$ \unit{\pound}.
Determine the steady state component of the displacement for $t>0$.

\item\label{exer:6.2.14}A \qty{32}{\pound} weight stretches a spring \qty{1}{\foot} in equilibrium. It is attached
to a dashpot with constant $c=\qty{12}{\pound\second\per\foot}$. The weight is initially
displaced 8 inches above equilibrium and released from rest. Find its
displacement for $t>0$.

\item\label{exer:6.2.15}A mass of one \unit{\kilo\gram} stretches a spring \qty{49}{\centi\meter} in equilibrium. A dashpot
attached to the spring supplies a damping force of \qty{4}{\newton} for each \unit{\meter\per\second}
of speed. The mass is initially displaced \qty{10}{\centi\meter} above equilibrium and
given a downward velocity of 1 m/sec. Find its displacement for $t>0$.

\item\label{exer:6.2.16}A mass of \qty{100}{\gram} stretches a spring \qty{98}{\centi\meter} in equilibrium. A dashpot
attached to the spring supplies a damping force of \qty{600}{\dyne} for each
\unit{\centi\meter\per\second} of speed. The mass is initially displaced \qty{10}{\centi\meter} above
equilibrium and given a downward velocity of \qty{1}{\meter\per\second}. Find its
displacement for $t>0$.

\item\label{exer:6.2.17}A \qty{192}{\pound} weight is suspended from a spring with constant $k=\qty{6}{\pound\per\foot}$
and subjected to an external force $F(t)=8\cos3t$ \unit{\pound}. Find the steady
state component of the displacement for $t>0$ if the medium resists
the motion with a force  equal to 8 times the speed in
\unit{\feet\per\second}.

\item\label{exer:6.2.18}A \qty{2}{\gram} mass is attached to a spring with constant \qty{20}{\dyne\per\centi\meter}. Find the
steady state component of the displacement if the mass is subjected to
an external force $F(t)=3\cos4t-5\sin4t$ dynes and a dashpot supplies
4 dynes of damping for each \unit{\centi\meter\per\second} of velocity.

\item\label{exer:6.2.19}\CGex
A \qty{96}{\pound} weight is attached to a spring with constant \qty{12}{\pound\per\foot}. Find
and graph the steady state component of the displacement if the mass is
subjected to an external force $F(t)=18\cos t-9\sin t$ \unit{\pound} and a
dashpot supplies \qty{24}{\pound} of damping for each \unit{\foot\per\second} of velocity.

\item\label{exer:6.2.20}A mass of one \unit{\kilo\gram} stretches a spring \qty{49}{\centi\meter} in equilibrium. It is
attached to a dashpot that supplies a damping force of \qty{4}{\newton} for each
\unit{\meter\per\second} of speed. Find the steady state component of its displacement if
it's subjected to an external force $F(t)=8\sin2t-6\cos2t$ \unit{\newton}.


\item\label{exer:6.2.21}A mass $m$ is suspended from a spring with constant $k$ and subjected to an
external force $F(t)=\alpha\cos\omega_0t+\beta\sin\omega_0t$, where
$\omega_0$ is the natural frequency of the spring-mass system without
damping. Find the steady state component of the displacement if a dashpot
with constant $c$ supplies damping.

\item\label{exer:6.2.22}Show that if $c_1$ and $c_2$ are not both zero then
\[
y=e^{r_1t}(c_1+c_2t)
\]
can't equal zero for more than one value of
$t$.

\item\label{exer:6.2.23}Show that if $c_1$ and $c_2$ are not both zero then
\[
y=c_1e^{r_1t}+c_2e^{r_2t}
\]
 can't equal zero for more than one
value of $t$.

\item\label{exer:6.2.24}Find the solution of the initial value problem
\[
my''+cy'+ky=0,\quad y(0)=y_0,\;y'(0)=v_0,
\]
given that the motion is underdamped, so  the general solution of
the equation is
\[
y=e^{-ct/2m}(c_1\cos\omega_1t+c_2\sin\omega_1t).
\]

\item\label{exer:6.2.25}Find the solution of the initial value problem
\[
my''+cy'+ky=0,\quad y(0)=y_0,\;y'(0)=v_0,
\]
given that the motion is overdamped, so  the general solution of
the equation is
\[
y=c_1e^{r_1t}+c_2e^{r_2t}\;(r_1,r_2<0).
\]

\item\label{exer:6.2.26}Find the solution of the initial value problem
\[
my''+cy'+ky=0,\quad y(0)=y_0,\;y'(0)=v_0,
\]
given that the motion is critically damped, so that the general
solution of the equation is of the form
\[
y=e^{r_1t}(c_1+c_2t)\,(r_1<0).
\]
\end{exerciselist}


\section{The \texorpdfstring{$RLC$}{RLC} Circuit}\label{sec:rlc}
\index{circuit!RLC}\index{RLC circuit}

In this section we consider the \emph{\color{blue}$RLC$ circuit}, shown
schematically in \zcref{figure:6.3.1}. As we'll see, the $RLC$
circuit is an electrical analog of a spring-mass system with
damping.

Nothing happens while the switch is open (dashed line). When the
switch is closed (solid line) we say that the \emph{\color{blue}circuit is
closed}\index{closed circuit}. Differences in electrical potential in a closed circuit
cause current to flow in the circuit. The battery or generator in
\zcref{figure:6.3.1} creates a difference in electrical potential
$E=E(t)$ between its two terminals, which we've marked arbitrarily as
positive and negative. (We could just as well interchange the
markings.) We'll say that $E(t)>0$ if the potential at the positive
terminal is greater than the potential at the negative terminal,
$E(t)<0$ if the potential at the positive terminal is less than the
potential at the negative terminal, and $E(t)=0$ if the potential is
the same at the two terminals. We call $E$ the \emph{\color{blue}impressed
voltage}.

\begin{figure}
  \centering\color{blue}
  \includegraphics[alt={A circuit with resistor (resistance R), induction coil (inductance L), capacitor (capacitance C), a switch, and a battery or generator supplying an impressed voltage E=E(t).},
  width=5.67in,height=3.3in,keepaspectratio]{fig060301}
\caption{An $RLC$ circuit}
  \label{figure:6.3.1}
\end{figure}


At any  time $t$, the same current flows in all points of the
circuit. We denote current by $I=I(t)$. We say that $I(t)>0$ if the
direction of flow is around the circuit from the positive terminal of
the battery or generator back to the negative terminal, as indicated
by the arrows in \zcref{figure:6.3.1}   $I(t)<0$ if the flow is in the
opposite direction, and $I(t)=0$ if no current flows at time $t$.

Differences in potential occur at the resistor, induction coil, and
capacitor in \zcref{figure:6.3.1}. Note that the two sides of each
of these components are also identified as positive and negative. The
\emph{\color{blue}voltage drop across}  each component is defined to be the
potential on the positive side of the component minus the potential on
the negative side. This terminology is somewhat misleading, since
``drop'' suggests a decrease even though changes in potential are
signed quantities and therefore may be increases. Nevertheless,
we'll go along with tradition and call them voltage drops.
The voltage drop across the resistor in \zcref{figure:6.3.1} is given
by
\begin{equation} \label{eq:6.3.1}
V_R=IR,
\end{equation}
where $I$ is current and $R$ is a positive constant,
the \indexterm{resistance} of the resistor.
The voltage drop across the induction coil is given by
\begin{equation}\label{eq:6.3.2}
V_I=L\frac{dI}{dt}=LI',
\end{equation}
where $L$ is a positive constant, the \indexterm{inductance} of the coil.

A \indexterm{capacitor} stores electrical charge $Q=Q(t)$, which is related to
the current in the circuit by the equation
\begin{equation} \label{eq:6.3.3}
Q(t)=Q_0+\int_0^tI(\tau)\,d\tau,
\end{equation}
where $Q_0$ is the \indexterm{charge} on the capacitor at $t=0$. The voltage drop
across a capacitor is given by
\begin{equation} \label{eq:6.3.4}
V_C=\frac{Q}{C},
\end{equation}
where $C$ is a positive constant, the \indexterm{capacitance} of the
capacitor.

\zcref[S]{table:6.3.1} names the units for the quantities that we've
discussed. The units are defined so that
\begin{align*}
1\text{\,volt}&=1\text{\,ampere}\cdot1\text{\,ohm}\\
&=1\text{\,henry}\cdot1\,\text{ampere}/\text{second}\\
&=1\text{\,coulomb}/\text{\,farad}\qquad\text{and}\\
1\text{\,ampere}&=1\text{\,coulomb}/\text{second}.
\end{align*}

\begin{table}
\caption{Electrical Units}
\label{table:6.3.1}
\centering
\tagpdfsetup{table/header-rows={1}}
\begin{tabular}{ccc}\toprule
Symbol&Name&Unit\\\midrule
$E$  &  Impressed Voltage & volt \unit\volt  \\
$I$ & Current & ampere (\unit\ampere)\\
$Q$ & Charge & coulomb (\unit\coulomb)\\
$R$ & Resistance& ohm (\(\Omega\))\\
% todo see https://github.com/josephwright/siunitx/issues/873
%$R$ & Resistance& ohm (\unit\ohm)\\
$L$ & Inductance& henry (\unit\henry)\\
$C$ & Capacitance&farad (\unit\farad)\\
\bottomrule
\end{tabular}
\end{table}



According to
\personHref{Kirchhoff}{Kirchoff's law}\index{Kirchoff's law}, the sum of
the voltage drops in a closed $RLC$ circuit equals the impressed
voltage. Therefore, from \eqref{eq:6.3.1}, \eqref{eq:6.3.2}, and \eqref{eq:6.3.4},
\begin{equation}\label{eq:6.3.5}
 LI'+RI+\frac{1}{C}Q=E(t).
\end{equation}
This equation contains two unknowns, the current $I$ in the
circuit and the charge $Q$ on the capacitor. However, \eqref{eq:6.3.3}
implies that $Q'=I$, so \eqref{eq:6.3.5} can be converted into the second
order equation
\begin{equation} \label{eq:6.3.6}
 LQ''+RQ'+\frac{1}{C}Q=E(t)
\end{equation}
in $Q$. To find the current flowing in an $RLC$ circuit, we solve
\eqref{eq:6.3.6} for $Q$ and then differentiate the solution to obtain
$I$.

In \zcref{sec:spring_i,sec:spring_ii} we encountered the
equation
\begin{equation} \label{eq:6.3.7}
my''+cy'+ky=F(t)
\end{equation}
in connection with spring-mass systems. Except for notation this
equation is the same as \eqref{eq:6.3.6}. The correspondence between
electrical and mechanical quantities connected with \eqref{eq:6.3.6} and
\eqref{eq:6.3.7} is shown in \zcref{table:6.3.2}.


\begin{table}
\caption{Electrical and Mechanical Units}
\label{table:6.3.2}
\centering
\tagpdfsetup{table/header-rows={1}}
\begin{tabular}{lclc}\toprule
\multicolumn{2}{c}{Electrical}&
\multicolumn{2}{c}{Mechanical}\\\cmidrule(r){1-2}\cmidrule(l){3-4}
charge& $Q$& displacement&$y$\\
current&$I$&velocity&$y'$\\
impressed voltage&$E(t)$&external force&$F(t)$\\
inductance&$L$&Mass&$m$\\
resistance&$R$&damping&$c$\\
1/capacitance&$1/C$&spring constant&$k$\\
\bottomrule
\end{tabular}
\end{table}

The equivalence between \eqref{eq:6.3.6} and \eqref{eq:6.3.7} is an example of
how mathematics unifies fundamental similarities in diverse physical
phenomena. Since we've already studied the properties of solutions of
\eqref{eq:6.3.7} in \zcref{sec:spring_i,sec:spring_ii}, we
can obtain
results concerning solutions of \eqref{eq:6.3.6} by simpling changing
notation, according to \zcref{table:6.3.1}.


\boxit{Free Oscillations}
\index{free oscillations}
\index{oscillation!free}

We say that an $RLC$ circuit is in \emph{\color{blue}free oscillation} if $E(t)=0$
for $t>0$, so that \eqref{eq:6.3.6} becomes
\begin{equation} \label{eq:6.3.8}
LQ''+RQ'+\frac{1}{C}Q=0.
\end{equation}
The characteristic equation of \eqref{eq:6.3.8} is
\[
Lr^2+Rr+\frac{1}{C}=0,
\]
with roots
\begin{equation}\label{eq:6.3.9}
r_1=\frac{-R-\sqrt{R^2-4L/C}}{2L}\quad\text{and}\quad  r_2=
\frac{-R+\sqrt{R^2-4L/C}}{2L}.
\end{equation}
There are three cases to consider, all analogous to the cases
considered in \zcref{sec:spring_ii} for free vibrations of a damped
spring-mass system.

\begin{enumext}[label=\arabic*,wrap-label={\textsc{Case} #1.},list-indent=0pt,labelwidth=3.5em]
\item
The oscillation is \emph{\color{blue}underdamped} if $R<\sqrt{4L/C}$. In this
case, $r_1$ and $r_2$ in \eqref{eq:6.3.9} are complex conjugates, which we
write as
\[
r_1=-\frac{R}{2L}+i\omega_1\quad\text{and}\quad  r_2=-\frac{R}{2L}-i\omega_1,
\]
where
\[
\omega_1=\frac{\sqrt{4L/C-R^2}}{2L}.
\]
The general solution of  \eqref{eq:6.3.8}  is
\[
Q=e^{-Rt/2L}(c_1\cos\omega_1 t+c_2\sin\omega_1 t),
\]
which we can write as
\begin{equation}\label{eq:6.3.10}
Q=Ae^{-Rt/2L}\cos(\omega_1 t-\phi),
\end{equation}
 where
\[
A=\sqrt{c_1^2+c_2^2},\quad A\cos\phi=c_1,\quad\text{and}\quad  A\sin\phi=c_2.
\]
In the idealized case where $R=0$, the solution \eqref{eq:6.3.10}
reduces to
\[
Q=A\cos\left(\frac{t}{\sqrt{LC}}-\phi\right),
\]
which is analogous to  the simple harmonic motion of an undamped
spring-mass system in free vibration.

Actual $RLC$ circuits are usually underdamped, so the case we've just
considered is the most important. However, for completeness we'll
consider the other two possibilities.

\item
The oscillation is \emph{\color{blue}overdamped} if $R>\sqrt{4L/C}$. In this case,
the zeros $r_1$ and $r_2$ of the characteristic polynomial are real,
with $r_1<r_2<0$ (see \eqref{eq:6.3.9}), and the general solution of
\eqref{eq:6.3.8} is
\begin{equation}\label{eq:6.3.11}
Q=c_1e^{r_1t}+c_2e^{r_2t}.
\end{equation}

\item
The oscillation is \emph{\color{blue}critically damped} if $R=\sqrt{4L/C}$. In
this case, $r_1=r_2=-R/2L$ and the general solution of \eqref{eq:6.3.8} is
\begin{equation}\label{eq:6.3.12}
Q=e^{-Rt/2L}(c_1+c_2t).
\end{equation}
\end{enumext}

If $R\ne0$, the exponentials in \eqref{eq:6.3.10}, \eqref{eq:6.3.11},
and
\eqref{eq:6.3.12} are negative, so the solution of any homogeneous initial
value problem
\[
 LQ''+RQ'+\frac{1}{C}Q=0,\quad Q(0)=Q_0,\quad Q'(0)=I_0,
\]
approaches zero exponentially as $t\to\infty$. Thus, all such
solutions are \emph{\color{blue}transient}\index{transient!solution}, in the sense defined
\zcref{sec:spring_ii} in the discussion of forced vibrations of a
spring-mass system with damping.

\begin{example}\label{example:6.3.1}
At $t=0$ a current of 2 amperes flows in an $RLC$ circuit with
resistance $R=40$ ohms, inductance $L=.2$ henrys, and capacitance
$C=10^{-5}$ farads. Find the current flowing in the circuit at $t>0$
if the initial charge on the capacitor is 1 coulomb. Assume that
$E(t)=0$ for $t>0$.
\end{example}

\begin{solution}
The equation for the charge $Q$ is \[
\frac{1}{5}Q''+40Q'+10000Q=0, \] or
\begin{equation} \label{eq:6.3.13}
Q''+200Q'+50000Q=0.
\end{equation}
Therefore we must solve the initial value problem
\begin{equation} \label{eq:6.3.14}
Q''+200Q'+50000Q=0,\quad Q(0)=1,\quad Q'(0)=2.
\end{equation}
The desired current is the derivative of the solution of this initial
value problem.

The characteristic equation of \eqref{eq:6.3.13} is
\[
r^2+200r+50000=0,
\]
which has complex zeros $r=-100\pm200i$. Therefore the general
solution of \eqref{eq:6.3.13} is
\begin{equation} \label{eq:6.3.15}
Q=e^{-100t}(c_1\cos200t+c_2\sin200t).
\end{equation}
Differentiating this and collecting like terms yields
\begin{equation} \label{eq:6.3.16}
Q'=-e^{-100t}\left[(100c_1-200c_2)\cos200t+
(100c_2+200c_1)\sin200t\right].
\end{equation}
To find the solution of the initial value problem \eqref{eq:6.3.14},
we set $t=0$ in \eqref{eq:6.3.15} and \eqref{eq:6.3.16} to obtain
\[
c_1=Q(0)=1\quad\text{and}\quad    -100c_1+200c_2=Q'(0)=2;
\]
therefore, $c_1=1$ and $c_2=51/100$, so
\[
Q=e^{-100t}\left(\cos200t+\frac{51}{100}\sin200t\right)
\]
is the solution of \eqref{eq:6.3.14}.
Differentiating this yields
\[
I=e^{-100t}(2\cos200t-251\sin200t).\qedhere
\]
\end{solution}


\boxit{Forced Oscillations With Damping}
\index{forced oscillation!with damping}
\index{oscillation!forced with damping}
\index{damping!forced oscillation with}

An initial value problem for \eqref{eq:6.3.6}  has the form
\begin{equation} \label{eq:6.3.17}
 LQ''+RQ'+\frac{1}{C}Q=E(t),\quad Q(0)=Q_0,\quad Q'(0)=I_0,
\end{equation}
where $Q_0$ is the initial charge on the capacitor and $I_0$ is the
initial current in the circuit. We've already seen that if $E\equiv0$
then all solutions of \eqref{eq:6.3.17} are transient. If
$E\not\equiv0$, we know that the solution of \eqref{eq:6.3.17} has
the form $Q=Q_c+Q_p$, where $Q_c$ satisfies the complementary
equation, and approaches zero exponentially as $t\to\infty$ for any
initial conditions , while $Q_p$ depends only on $E$ and is
independent of the initial conditions. As in the case of forced
oscillations of a spring-mass system with damping, we call $Q_p$ the
\emph{\color{blue}steady state}\index{charge!steady state} charge on the capacitor of the $RLC$ circuit. Since
$I=Q'=Q_c'+Q_p'$ and $Q_c'$ also tends to zero exponentially as
$t\to\infty$, we say that $I_c=Q'_c$ is the \emph{\color{blue}transient current}\index{current!transient}\index{transient!current}
and $I_p=Q_p'$ is the \emph{\color{blue}steady state} current\index{steady state!current}\index{current!steady state}. In most
applications we're interested only in the steady state charge and
current.

\begin{example}\label{example:6.3.2}
Find the amplitude-phase form of the steady state current in the $RLC$
circuit in \zcref{figure:6.3.1} if the impressed voltage, provided by
an alternating current generator, is $E(t)=E_0\cos\omega t$.
\end{example}

\begin{solution}
We'll first find the steady state charge on the capacitor as
a particular solution of
\[
LQ''+RQ'+\frac{1}{C}Q=E_0\cos\omega t.
\]
To do, this we'll simply reinterpret a result obtained in
\zcref{sec:spring_ii}, where we found that the steady state solution of
\[
my''+cy'+ky=F_0\cos\omega t
\]
is
\[
y_p=\frac{F_0}{\sqrt{(k-m\omega^2)^2+c^2\omega^2}}
\cos(\omega t-\phi),
\]
where
\[
\cos\phi=\frac{k-m\omega^2}{\sqrt
{(k-m\omega^2)^2+c^2\omega^2}}\quad\text{and}\quad 
\sin\phi=\frac{c\omega}{\sqrt{(k-m\omega^2)^2+c^2\omega^2}}.
\]
(See \zcref{eq:6.2.14,eq:6.2.15}.) By making the
appropriate changes in the symbols
(according to \zcref{table:6.3.2}) yields the steady state charge
\[
Q_p=\frac{E_0}{\sqrt{(1/C-L\omega^2)^2+R^2\omega^2}}\cos(\omega
t-\phi),
\]
where
\[
\cos\phi=\frac{1/C-L\omega^2}{\sqrt{(1/C-L\omega^2)^2+R^2\omega^2}}
\quad\text{and}\quad 
\sin\phi=\frac{R\omega}{\sqrt{(1/C-L\omega^2)^2+R^2\omega^2}}.
\]
 Therefore the steady state current in the circuit
is
\[
I_p=Q_p'=
-\frac{\omega E_0}{\sqrt{(1/C-L\omega^2)^2+R^2\omega^2}}\sin(\omega
t-\phi).\qedhere
\]
\end{solution}

\exercises
In \zcref[range]{exer:6.3.1,exer:6.3.5} find the current in the
$RLC$ circuit, assuming that $E(t)=0$ for $t>0$.

\begin{exerciselist}
\item\label{exer:6.3.1}$R=3$ ohms;  \, $L=.1$ henrys;  \, $C=.01$ farads;   $Q_0=0$
coulombs;  \,$I_0=2$ amperes.

\item\label{exer:6.3.2}$R=2$ ohms;  \, $L=.05$ henrys;  \, $C=.01$ farads;  $Q_0=2$
coulombs;  \,$I_0=-2$ amperes.

\item\label{exer:6.3.3}$R=2$ ohms;  \, $L=.1$ henrys;  \, $C=.01$ farads;   $Q_0=2$
coulombs;  \,$I_0=0$ amperes.

\item\label{exer:6.3.4}$R=6$ ohms;  \, $L=.1$ henrys;  \, $C=.004$ farads;  $Q_0=3$
coulombs;  \,$I_0=-10$ amperes.

\item\label{exer:6.3.5}$R=4$ ohms;  \, $L=.05$ henrys;  \, $C=.008$ farads;   $Q_0=-1$
coulombs;  \,$I_0=2$ amperes.


\exercisetext{In \zcref[range]{exer:6.3.6,exer:6.3.10} find the steady
state current in the circuit described by the  equation.}

\item\label{exer:6.3.6}$\dfrac{1}{10}Q''+3Q'+100Q=5\cos10t-5\sin10t$

\item\label{exer:6.3.7}$\dfrac{1}{20}Q''+2Q'+100Q=10\cos25t-5\sin25t$

\item\label{exer:6.3.8}$\dfrac{1}{10}Q''+2Q'+100Q=3\cos50t-6\sin50t$

\item\label{exer:6.3.9}$\dfrac{1}{10}Q''+6Q'+250Q=10\cos100t+30\sin100t$

\item\label{exer:6.3.10}$\dfrac{1}{20}Q''+4Q'+125Q=15\cos30t-30\sin30t$

\item\label{exer:6.3.11}Show that if $E(t)=U\cos\omega t+V\sin\omega t$ where $U$
and $V$ are constants then the steady state current in
the $RLC$ circuit shown in \zcref{figure:6.3.1} is
\[
I_p=\frac{\omega^2RE(t)+(1/C-L\omega^2)E'(t)}{\Delta},
\]
where
\[
\Delta=(1/C-L\omega^2)^2+R^2\omega^2.
\]

\item\label{exer:6.3.12}Find the amplitude of the steady state current $I_p$ in the $RLC$
circuit shown in \zcref{figure:6.3.1}  if $E(t)=U\cos\omega
t+V\sin\omega t$, where $U$ and $V$ are constants. Then find the value
$\omega_0$ of $\omega$ maximizes the amplitude, and
 find the maximum amplitude.

\exercisetext{In \zcref[range]{exer:6.3.13,exer:6.3.17} plot the
amplitude of the steady state current against $\omega$. Estimate the
value of $\omega$ that maximizes the amplitude of the steady state
current, and estimate this maximum amplitude. \hint{You can confirm
your results by doing \zcref{exer:6.3.12}.}}

\item\label{exer:6.3.13}\Lex
$\dfrac{1}{10}Q''+3Q'+100Q=U\cos\omega t+V\sin\omega t$

\item\label{exer:6.3.14}\Lex
$\dfrac{1}{20}Q''+2Q'+100Q=U\cos\omega t+V\sin\omega t$

\item\label{exer:6.3.15}\Lex
$\dfrac{1}{10}Q''+2Q'+100Q=U\cos\omega t+V\sin\omega t$

\item\label{exer:6.3.16}\Lex
$\dfrac{1}{10}Q''+6Q'+250Q=U\cos\omega t+V\sin\omega t$

\item\label{exer:6.3.17}\Lex
$\dfrac{1}{20}Q''+4Q'+125Q=U\cos\omega t+V\sin\omega t$
\end{exerciselist}


\section{Motion Under a Central Force}\label{sec:central_force}
\index{motion!under a central force}
\index{central force!motion under}

We'll now study the motion of a object moving under the influence
of a \emph{\color{blue}central force}; that is, a force whose magnitude at any
point $P$ other than the origin depends only on the distance from
$P$ to the origin, and whose direction at $P$ is parallel to the line
connecting $P$ and the origin, as indicated in
\zcref{figure:6.4.1} for the case where the direction of the force
at every point is toward the origin.
 Gravitational forces are central
forces;   for example, as mentioned in \zcref{sec:elementary_mechanics}, if we
assume that Earth is a perfect sphere with constant mass density then
Newton's law of gravitation\index{Newton's!law of gravitation} asserts that the force exerted on an
object by Earth's gravitational field is proportional to the mass of
the object and inversely proportional to the square of its distance
from the center of Earth, which we take to be the origin.

If the initial position and velocity vectors of an object moving under
a central force are parallel, then the subsequent motion is along the
line from the origin to the initial position. Here we'll assume that
the initial position and velocity vectors are not parallel;   in this
case the subsequent motion is in the plane determined by them. For
convenience we take this to be the $xy$-plane. We'll
consider
the problem of determining the curve traversed by the object. We call
this curve the \emph{\color{blue}orbit}.

\begin{figure}
  \centering\color{blue}
  \includegraphics[alt={A vector field where everything points to the origin.},
  width=5.67in,height=3.3in,keepaspectratio]{fig060401}
  \caption{}
  \label{figure:6.4.1}
\end{figure}

We can represent a central force in terms of polar coordinates
\[
x=r\cos\theta,\quad y=r\sin\theta
\]
 as
\[
\vect{F}(r,\theta)=f(r)(\cos\theta\ivect+\sin\theta\jvect).
\]
We assume that $f$ is continuous for all $r>0$. The magnitude of $\vect{F}$ at $(x,y)=(r\cos\theta,r\sin\theta)$ is $|f(r)|$, so it depends
only on the distance $r$ from the point to the origin   the direction
of $\vect{F}$ is from the point to the origin if $f(r)<0$, or from the
origin to the point if $f(r)>0$. We'll show that the orbit of an
object with mass $m$ moving under this force is given by
\[
r(\theta)=\frac{1}{u(\theta)},
\]
where $u$  is solution of the differential equation
\begin{equation}\label{eq:6.4.1}
 \frac{d^2u}{d\theta^2}+u=-\frac{1}{mh^2u^2}f(1/u),
\end{equation}
and $h$ is a constant defined below.



Newton's second law of motion\index{Newton's!second law of motion} ($\vect{F}=m\vect{a}$) says that the polar
coordinates
$r=r(t)$ and $\theta=\theta(t)$ of
the particle  satisfy the vector differential equation
\begin{equation} \label{eq:6.4.2}
m(r\cos\theta\ivect+r\sin\theta\jvect)''=
f(r)(\cos\theta\ivect+\sin\theta\jvect).
\end{equation}
To deal with this equation we introduce  the unit vectors
\[
\vect{e}_1=\cos\theta\ivect+\sin\theta\jvect
\quad\text{and}\quad 
\vect{e}_2=-\sin\theta\ivect+\cos\theta\jvect.
\]
Note that   $\vect{e}_1$ points in the
direction of increasing $r$ and $\vect{e}_2$ points in the
direction of increasing $\theta$ (\zcref{figure:6.4.2});
moreover,
\begin{equation} \label{eq:6.4.3}
\frac{d\vect{e}_1}{d\theta}=\vect{e}_2,\quad
\frac{d\vect{e}_2}{d\theta}=-\vect{e}_1,
\end{equation}
and
\[
\vect{e}_1\cdot\vect{e}_2=\cos\theta(-\sin\theta)+\sin\theta\cos\theta=0,
\]
so $\vect{e}_1$ and $\vect{e}_2$ are  perpendicular.
Recalling that the single prime $(')$ stands for differentiation with
respect to
$t$, we see from \eqref{eq:6.4.3}  and the chain rule that
\begin{equation} \label{eq:6.4.4}
\vect{e}_1'=\theta'\vect{e}_2\quad\text{and}\quad 
\vect{e}_2'=-\theta'\vect{e}_1.
\end{equation}

\begin{figure}
  \centering\color{blue}
  \includegraphics[alt={Two vectors pointing away from the origin.  Vector e1 points in the (1,1) direction.  Vector e2 points in the (-1,1) direction.},
  width=5.67in,height=3.3in,keepaspectratio]{fig060402}
  \caption{}
  \label{figure:6.4.2}
\end{figure}



Now we can write \eqref{eq:6.4.2} as
\begin{equation} \label{eq:6.4.5}
m(r\vect{e}_1)''=f(r)\vect{e}_1.
\end{equation}
But
\[
(r\vect{e}_1)'=r'\vect{e}_1+r\vect{e}_1'=r'\vect{e}_1+r\theta'\vect{e}_2
\]
(from \eqref{eq:6.4.4}), and
\begin{align*}
(r\vect{e}_1)''&=(r'\vect{e}_1+r\theta'\vect{e}_2)'\\
&=r''\vect{e}_1+r'\vect{e}_1'+(r\theta''+r'\theta')\vect{e}_2+r\theta'\vect{e}_2'\\
&=r''\vect{e}_1+r'\theta'\vect{e}_2
+(r\theta''+r'\theta')\vect{e}_2 -r(\theta')^2\vect{e}_1\quad\text{(from \eqref{eq:6.4.4})}\\
&=\left(r''-r(\theta')^2\right)\vect{e}_1+(r\theta''+2r'\theta')\vect{e}_2.
\end{align*}
Substituting this into \eqref{eq:6.4.5} yields
\[
m\left(r''-r(\theta')^2\right)\vect{e}_1+m(r\theta''+2r'\theta')\vect{e}_2=f(r)\vect{e}_1.
\]
By equating the coefficients of $\vect{e}_1$ and $\vect{e}_2$ on the
two sides of this equation we see that
\begin{equation} \label{eq:6.4.6}
m\left(r''-r(\theta')^2\right)=f(r)
\end{equation}
and
\[
r\theta''+2r'\theta'=0.
\]
Multiplying the last equation by $r$ yields
\[
r^2\theta''+2rr'\theta'=(r^2\theta')'=0,
\]
so
\begin{equation} \label{eq:6.4.7}
r^2\theta'=h,
\end{equation}
where $h$ is a constant that we can write in terms of the initial
conditions as
\[
h=r^2(0)\theta'(0).
\]
Since the initial position and velocity vectors are
\[
 r(0)\vect{e}_1(0)\quad\text{and}\quad 
r'(0)\vect{e}_1(0)+r(0)\theta'(0)\vect{e}_2(0),
\]
our assumption that these two vectors are not parallel implies that
$\theta'(0)\ne0$, so $h\ne0$.


Now let $u=1/r$.
Then  $u^2=\theta'/h$ (from \eqref{eq:6.4.7}) and
\[
r'=-\dfrac{u'}{u^2}=-h\left(\frac{u'}{\theta'}\right),
\]
which implies that
\begin{equation} \label{eq:6.4.8}
r'=-h\frac{du}{d\theta},
\end{equation}
since
\[
\frac{u'}{\theta'}=\frac{du}{dt}\left/\frac{d\theta}{dt}\right.=
\frac{du}{d\theta}.
\]
Differentiating \eqref{eq:6.4.8} with respect to $t$ yields
\[
r''=-h\frac{d}{dt}\left(\frac{du}{d\theta}\right)
=-h\frac{d^2u}{d\theta^2}\theta',
\]
which implies that
\[
r''=-h^2u^2\frac{d^2u}{d\theta^2}
\quad\text{since}\quad\theta'=hu^2.
\]
Substituting from these equalities into
\eqref{eq:6.4.6} and recalling that $r=1/u$ yields
\[
-m\left(h^2u^2\frac{d^2u}{d\theta^2}+\frac{1}{u}h^2u^4\right)
=f(1/u),
\]
and dividing through by $-mh^2u^2$ yields \eqref{eq:6.4.1}.


\zcref[S]{eq:6.4.7} has the following geometrical interpretation, which
is known as
\personHref{Kepler}{Kepler's second law}\index{Kepler's second law}.

\begin{theorem}\color{blue}\label{thmtype:6.4.1}
 The position vector
 of an object moving under a central force
sweeps out equal areas in equal times; more precisely, if
$\theta(t_1)\le\theta(t_2)$ then the (signed) area of the sector
\[
\{(x,y)=(r\cos\theta,r\sin\theta)\;  :\;  0\le r\le r(\theta),\;
\theta(t_1)\le\theta(t_2)\}
\]
(\zcref{figure:6.4.3}) is given by
\[
A=\frac{h(t_2-t_1)}{2},
\]
where $h=r^2\theta'$, which we have shown to be constant.
\end{theorem}

\begin{figure}
  \centering\color{blue}
  \includegraphics[alt={A circular sector in the first quadrant with its vertex at the origin and between the angles 𝜃=𝜃(t1) and 𝜃=𝜃(t2).},
  width=5.67in,height=3.3in,keepaspectratio]{fig060403}
  \caption{}
  \label{figure:6.4.3}
\end{figure}

\begin{proof}
Recall from calculus that the area of the shaded sector
in \zcref{figure:6.4.3} is
\[
A=\frac{1}{2}\int_{\theta(t_1)}^{\theta(t_2)} r^2(\theta)\,d\theta,
\]
where $r=r(\theta)$ is the polar representation of the orbit.
Making the change of variable $\theta=\theta(t)$ yields
\begin{equation} \label{eq:6.4.9}
A=\frac{1}{2}\int_{t_1}^{t_2} r^2(\theta(t))\theta'(t)\,dt.
\end{equation}
 But
\eqref{eq:6.4.7}  and \eqref{eq:6.4.9} imply that
\[
A=\frac{1}{2}\int_{t_1}^{t_2} h\,dt=\frac{h(t_2-t_1)}{2},
\]
which completes the proof.
\end{proof}


\boxit{Motion Under an Inverse Square Law Force}
\index{motion!with inverse square law force}
\index{inverse square law force!motion with}

In the special case where
 $f(r)=-mk/r^2=-mku^2$, so $\vect{F}$  can be interpreted
as a gravitational force, \eqref{eq:6.4.1}  becomes
\begin{equation} \label{eq:6.4.10}
 \frac{d^2u}{d\theta^2}+u=\frac{k}{h^2}.
\end{equation}
The general solution of the complementary equation
\[
 \frac{d^2u}{d\theta^2}+u=0
\]
can be written in amplitude-phase form as
\[
u=A\cos(\theta-\phi),
\]
where $A\ge0$ and $\phi$ is a phase angle.
Since $u_p=k/h^2$ is a particular solution of \eqref{eq:6.4.10},
the general solution of \eqref{eq:6.4.10} is
\[
u=A\cos(\theta-\phi)+\frac{k}{h^2};
\]
hence, the orbit is given by
\[
r=\left(A\cos(\theta-\phi)+\frac{k}{h^2}\right)^{-1},
\]
which we rewrite as
\begin{equation} \label{eq:6.4.11}
r=\frac{\rho}{1+e\cos(\theta-\phi)},
\end{equation}
where
\[
\rho=\frac{h^2}{k}\quad\text{and}\quad  e=A\rho.
\]
A curve satisfying \eqref{eq:6.4.11} is a conic section with a focus at
the origin (\zcref{exer:6.4.1}). The nonnegative constant $e$ is the
\emph{\color{blue}eccentricity}\index{eccentricity of orbit}\index{orbit!eccentricity}
of the orbit, which is
an ellipse\index{elliptic orbit}\index{orbit!elliptic} if
$e<1$
(a circle\index{circular orbit}\index{orbit!circular} if $e=0$),
a parabola if $e=1$, or a hyperbola if $e>1$.

\begin{figure}
  \centering\color{blue}
  \includegraphics[alt={An ellipse with one focus at the origin.  The point closest to the origin is labeled P; the point furthest from the origin is labeled A.  The ellipse is not aligned to either axis, so that A is in the first quadrant and P is in the third quadrant.},
  width=5.67in,height=3.3in,keepaspectratio]{fig060404}
  \caption{}
  \label{figure:6.4.4}
\end{figure}

If  the orbit is an ellipse, then the minimum and maximum
values of $r$  are
\begin{align*}
r_{\min}&=\frac{\rho}{1+e}\quad\text{(the \indexterm{perihelion
distance}, attained when
$\theta=\phi$)}  \\
r_{\max}&=\frac{\rho}{1-e}\quad\text{(the
\indexterm{aphelion distance}, attained when
$\theta=\phi+\pi$)}.
\end{align*}
\zcref[S]{figure:6.4.4} shows
a typical elliptic orbit.
The point $P$ on the orbit where $r=r_{\min}$ is the \indexterm{perigee} 
and the point $A$ where $r=r_{\max}$ is the \indexterm{apogee}.



For example, Earth's orbit around the Sun is approximately an ellipse
with $e\approx.017$, $r_{\min}\approx91\times 10^6$ miles, and
$r_{\max}\approx 95\times 10^6$ miles. Halley's comet has a very
elongated approximately elliptical orbit around the sun, with
$e\approx.967$, $r_{\min}\approx55\times10^6$ miles, and
$r_{\max}\approx33\times10^8$ miles. Some comets (the nonrecurring
type) have parabolic or hyperbolic orbits.

\exercises
\begin{exerciselist}


\item\label{exer:6.4.1}Find the equation of the curve
\[
r=\frac{\rho}{1+e\cos(\theta-\phi)}
\tag{A}
\]
in terms of
$(X,Y)=\left(r\cos(\theta-\phi),r\sin(\theta-\phi)\right)$,
which are rectangular coordinates with respect to the axes shown in
\zcref{figure:6.4.5}. Use your results to verify that
(A)
is the equation of an ellipse if $0<e<1$, a parabola if $e=1$,
or a hyperbola if $e>1$. If $e<1$, leave your answer in
the form
\[
\frac{(X-X_0)^2}{a^2}+\frac{(Y-Y_0)^2}{b^2}=1,
\]
and show that  the area of the ellipse is
\[
A=\frac{\pi\rho^2}{(1-e^2)^{3/2}}.
\]
Then use \zcref{thmtype:6.4.1} to show that the time required for
the object to traverse the entire orbit is
\[
T=\frac{2\pi\rho^2}{h(1-e^2)^{3/2}}.
\]
(This is
\personHref{Kepler}{Kepler's third law}\index{Kepler's third law};
$T$ is called the \emph{\color{blue}period}\index{period of orbit}\index{orbit!period of} of the orbit.)

\begin{minipage}{\linewidth}
\tagstructbegin{tag=Sect}
  \centering
  \includegraphics[alt={An ellipse with one focus at the origin.  The ellipse is not aligned to either axis, so that the aphelion is in the third quadrant and the perihelion is in the first quadrant.  An axis X is along the major diameter; an axis Y is along the minor diameter.},
  width=5.67in,height=3.3in,keepaspectratio]{fig060405}
\captionof{figure}{}
  \label{figure:6.4.5}
\tagstructend
\end{minipage}


\item\label{exer:6.4.2}Suppose an object with mass $m$ moves in the $xy$-plane under the
central force
\[
\vect{F}(r,\theta)=-\frac{mk}{r^2}(\cos\theta\ivect+\sin\theta\jvect),
\]
where $k$ is a positive constant. As we shown, the orbit of the
object is given by
\[
r=\frac{\rho}{1+e\cos(\theta-\phi)}.
\]
Determine $\rho$, $e$, and $\phi$ in terms of the initial conditions
\[
r(0)=r_0,\quad  r'(0)=r_0',\qquad\text{and}\qquad\theta(0)=\theta_0,\quad
\theta'(0)=\theta_0'.
\]
Assume that the initial position and velocity vectors are not
collinear.


\item\label{exer:6.4.3}Suppose we wish to put a satellite with mass $m$ into an
elliptical orbit around Earth. Assume that the only force acting on
the object is Earth's gravity, given by
\[
\vect{F}(r,\theta)=-mg\left(\frac{R^2}{r^2}\right)(\cos\theta\ivect+\sin\theta\jvect),
\]
where $R$ is Earth's radius, $g$ is the acceleration due to gravity at
Earth's surface, and $r$ and $\theta$ are polar coordinates in the
plane of the orbit, with the origin at Earth's center.

\begin{alist}
\item % (a)
Find the eccentricity required to make the aphelion and perihelion
distances equal to $R\gamma_1$ and $R\gamma_2$, respectively, where
$1<\gamma_1<\gamma_2$.
\item % (b)
Find the  initial conditions
\[
r(0)=r_0,\quad  r'(0)=r_0', \text{\; and \;} \theta(0)=\theta_0,\quad
\theta'(0)=\theta_0'
\]
required to make the initial point the perigee, and the motion along
the orbit in the direction of increasing $\theta$. \hint{Use the
results of \zcref{exer:6.4.2}.}
\end{alist}


\item\label{exer:6.4.4}An object with mass $m$  moves in a spiral orbit
$r=c\theta^2$ under a central force
\[
\vect{F}(r,\theta)=f(r)(\cos\theta\ivect+\sin\theta\jvect).
\]
Find $f$.


\item\label{exer:6.4.5}An object with mass $m$  moves in the orbit
$r=r_0e^{\gamma\theta}$ under a central force
\[
\vect{F}(r,\theta)=f(r)(\cos\theta\ivect+\sin\theta\jvect).
\]
 Find $f$.


\item\label{exer:6.4.6}Suppose an object with mass $m$ moves under the central force
\[
\vect{F}(r,\theta)=-\frac{mk}{r^3}(\cos\theta\ivect+\sin\theta\jvect),
\]
with
\[
r(0)=r_0,\quad  r'(0)=r_0', \text{\; and \;} \theta(0)=\theta_0,\quad
\theta'(0)=\theta_0',
\]
 where $h=r_0^2\theta_0'\ne0$.
\begin{alist}
\item % (a)
 Set up a second order initial
value problem for $u=1/r$ as a function of~$\theta$.
\item % (b)
Determine $r=r(\theta)$ if (i) $h^2<k$;\;   (ii) $h^2=k$;\;   (iii)
$h^2>k$.
\end{alist}

\end{exerciselist}
